\section{CAFs as Mediators of Treatment Resistance}
\label{sec:resistance}

A growing body of evidence indicates that CAFs contribute to therapeutic resistance in breast cancer through multiple mechanisms. Understanding these mechanisms is essential for developing strategies to overcome resistance and improve treatment outcomes.

\subsection{Chemoresistance}

CAFs promote resistance to chemotherapy through several interconnected pathways:

\subsubsection{Physical barrier and drug delivery}

The dense, collagen-rich ECM produced by myCAFs creates a physical barrier that impedes drug penetration into tumour tissue \citep{provenzano2012}. This barrier increases interstitial fluid pressure, compresses blood vessels, and reduces perfusion, limiting the delivery of systemically administered chemotherapeutics.

In breast cancer, high stromal density---measured by collagen content or imaging features---is associated with poor response to neoadjuvant chemotherapy \citep{farinha2019}. Preclinical studies have shown that strategies to reduce stromal density (e.g., anti-TGF-$\beta$ therapy, ATRA) can enhance drug delivery and improve chemosensitivity \citep{liu2023}.

\subsubsection{Paracrine survival signalling}

CAFs secrete multiple factors that directly promote tumour cell survival in the presence of chemotherapy:

\textbf{IL-6:} CAF-derived IL-6 activates STAT3 signalling in tumour cells, upregulating anti-apoptotic proteins (BCL-2, MCL-1) and drug efflux pumps (ABCB1, ABCG2) \citep{chang2014}. This paracrine survival signalling is a major mechanism of resistance to taxanes and anthracyclines in breast cancer.

\textbf{CXCL12:} CAF-secreted CXCL12 activates CXCR4 on tumour cells, triggering PI3K/AKT signalling that promotes survival and proliferation \citep{orimo2005stromal}. Inhibition of CXCR4 with plerixafor sensitizes breast cancer cells to chemotherapy in preclinical models.

\textbf{HGF:} Hepatocyte growth factor (HGF) secreted by CAFs activates MET signalling in tumour cells, promoting survival, invasion, and drug resistance \citep{wilson2014}. HGF/MET signalling is particularly relevant in basal-like breast cancer, where it contributes to resistance to EGFR-targeted therapies.

\subsubsection{Exosome-mediated resistance}

CAFs release exosomes containing miRNAs and proteins that confer chemoresistance to tumour cells \citep{zhao2023}:

\textbf{miR-21:} CAF-derived exosomes carry miR-21, which downregulates the pro-apoptotic protein PDCD4, enhancing tumour cell survival during chemotherapy \citep{zhao2023}.

\textbf{miR-200 family:} Paradoxically, CAF-derived exosomes can also carry miR-200 family members that suppress EMT and enhance chemosensitivity, highlighting the complex and context-dependent nature of exosome-mediated communication.

\textbf{Drug efflux pumps:} CAF-derived exosomes can transfer functional drug efflux pumps (P-glycoprotein) to tumour cells, directly conferring multidrug resistance \citep{zhao2023}.

\subsection{Endocrine resistance}

In ER+ breast cancer, CAFs contribute to resistance to endocrine therapies (tamoxifen, aromatase inhibitors) through multiple mechanisms:

\subsubsection{Alternative growth factor signalling}

CAFs secrete growth factors (FGF, HGF, IGF) that activate signalling pathways parallel to ER, providing tumour cells with ER-independent survival signals \citep{seth2021}. These alternative pathways bypass the requirement for estrogen signalling, rendering endocrine therapy ineffective.

In particular, CAF-derived FGF2 activates FGFR signalling in tumour cells, promoting proliferation and survival in the absence of estrogen \citep{seth2021}. FGFR inhibitors (erdafitinib, pemigatinib) are being evaluated in combination with endocrine therapy for ER+ breast cancer with CAF-mediated resistance.

\subsubsection{ECM-mediated resistance}

The ECM produced by CAFs provides physical and biochemical support that enhances tumour cell survival during endocrine therapy. Integrin signalling from ECM components activates FAK and PI3K/AKT pathways, promoting resistance to tamoxifen-induced apoptosis \citep{seth2021}.

\subsection{Immunotherapy resistance}

CAFs are major contributors to resistance to immune checkpoint inhibitors in breast cancer:

\subsubsection{Immune exclusion}

As discussed in Section~\ref{sec:immune}, myCAFs create a physical barrier that excludes immune cells from tumour nests, preventing the anti-tumour immune response required for ICI efficacy \citep{chhabra2023}. This immune-excluded phenotype is particularly common in TNBC with dense stroma, explaining the poor response to anti-PD-1/PD-L1 therapy in these tumours.

\subsubsection{Immunosuppressive milieu}

iCAFs create an immunosuppressive microenvironment through secretion of IL-6, CXCL12, and CCL2, which recruit and activate immunosuppressive cells (Tregs, MDSCs, M2 macrophages) \citep{erve2024}. This immunosuppressive milieu impairs T cell function and limits the efficacy of ICIs.

\subsubsection{Antigen presentation without co-stimulation}

apCAFs present tumour antigens to CD4+ T cells via MHC class II without adequate co-stimulation, inducing T cell anergy and tolerance \citep{erve2024}. This tolerogenic mechanism may contribute to the failure of anti-tumour immune responses despite ICI therapy.

\subsection{Targeted therapy resistance}

CAFs also contribute to resistance to targeted therapies in breast cancer:

\textbf{HER2+ breast cancer:} CAFs secrete neuregulin (NRG1) and other HER3 ligands that activate HER2/HER3 heterodimers, providing an alternative survival pathway when HER2 is blocked by trastuzumab or lapatinib \citep{seth2021}.

\textbf{CDK4/6 inhibitor resistance:} CAF-derived FGF activates FGFR signalling, which can bypass CDK4/6 inhibition by activating alternative cell cycle regulators \citep{seth2021}.

\textbf{PI3K inhibitor resistance:} CAFs secrete insulin-like growth factor (IGF) that activates PI3K signalling through IGF1R, providing an alternative activation pathway when PI3K is inhibited \citep{seth2021}.

\subsection{Clinical implications of CAF-mediated resistance}

Understanding CAF-mediated resistance has important clinical implications:

\textbf{Biomarker development:} Stromal features---including CAF density, collagen organization, and CAF subpopulation composition---may serve as predictive biomarkers for treatment response. High myCAF density, for example, may predict poor response to immunotherapy.

\textbf{Combination therapy:} Strategies that target CAFs in combination with conventional therapies may overcome resistance. For example, combining anti-TGF-$\beta$ therapy with anti-PD-1 may convert immune-excluded tumours into immune-inflamed tumours, enhancing ICI efficacy.

\textbf{Sequential therapy:} Understanding the temporal dynamics of CAF-mediated resistance may guide the sequencing of therapies. For example, CAF-targeted therapy prior to chemotherapy may enhance drug delivery and improve response.
